\section*{Bảng thuật ngữ}
\addcontentsline{toc}{section}{Bảng thuật ngữ}

\begin{spacing}{1.15}
\noindent
\small
\begin{tabularx}{\linewidth}{l L}
\toprule
\textbf{Thuật ngữ / Ký hiệu} & \textbf{Ý nghĩa / Định nghĩa} \\
\midrule
\texttt{API} & Application Programming Interface (Giao diện lập trình ứng dụng) \\
\texttt{AQE} & Adaptive Query Execution (Cơ chế tự động tối ưu hóa kế hoạch truy vấn khi đang chạy của Spark) \\
\texttt{CSV} & Comma-Separated Values (Định dạng file văn bản trong đó các giá trị được phân tách bằng dấu phẩy) \\
\texttt{DataFrame} & Cấu trúc dữ liệu dạng bảng được tối ưu hóa trong Spark Structured APIs \\
\texttt{GK} & Greenwald-Khanna (Thuật toán nén dữ liệu dùng để ước lượng các giá trị trung vị hoặc percentile một cách gần đúng) \\
\texttt{MapReduce} & Mô hình lập trình phân tán gồm hai pha chính: Map (ánh xạ dữ liệu) và Reduce (gom nhóm/thu gọn dữ liệu) \\
\texttt{Percentile} & Phân vị (ngưỡng giá trị chia tập dữ liệu thành các phần có tỷ lệ phần trăm tương ứng) \\
\texttt{RDD} & Resilient Distributed Dataset (Tập dữ liệu phân tán có khả năng phục hồi - mức trừu tượng cơ bản nhất của Spark) \\
\texttt{SKU} & Stock Keeping Unit (Mã định danh duy nhất cho từng mặt hàng cụ thể được lưu kho) \\
\texttt{Spark Structured API} & Các API cấp cao của Spark như DataFrame/Dataset và Spark SQL giúp tự động tối ưu hóa truy vấn \\
\texttt{StdDev (stddev\_pop)} & Population Standard Deviation (Độ lệch chuẩn tổng thể của một tập dữ liệu) \\
\texttt{Style} & Mã kiểu dáng của sản phẩm (một style có thể có nhiều SKU ứng với các kích cỡ khác nhau) \\
\texttt{Sliding Window} & Cửa sổ trượt (Kỹ thuật tính toán trên một khoảng thời gian động có độ dài cố định dịch chuyển liên tục) \\
\texttt{UDF} & User Defined Function (Hàm tự định nghĩa bởi lập trình viên để xử lý các logic phức tạp trong Spark) \\
\texttt{Variety} & Số lượng SKU phân biệt (distinct) thuộc về cùng một style của sản phẩm \\
\bottomrule
\end{tabularx}
\end{spacing}
